%======================================================================
\section{The ladder identity: the bias is carried by the squares with
no flippable ladder pair}\label{sec:ladder}

Session 13, after the component enumeration of \S\ref{sec:s13comp7}.
This section proves an exact identity that replaces the whole
chain \S\ref{sec:spectralexact}--\S\ref{sec:s13comp7} (pool, master
formula, annealed bias, transfer, trade chain) by a single
combinatorial statement.  Notation as in \S\ref{sec:s13trapped}:
$(L,P)\in X=X(n,\lambda;\alpha,\beta)$, $P=\{j,j'\}$, frame
$\pi=\pi_P$ ($L(\pi(r),j')=L(r,j)$), $C_r$ the $\pi$-cycle of row $r$,
$\rho_{x,y}$ the column permutation of rows $x,y$
($L(y,\rho_{x,y}(q))=L(x,q)$), $(x,y)$ \emph{flippable} iff $j,j'$ lie
in different cycles of $\rho_{x,y}$, and the \emph{$j'$-trade} at
$(x,y)$ exchanges rows $x,y$ on the $\rho_{x,y}$-cycle through $j'$,
which replaces $\pi$ by $(x\,y)\circ\pi$ if $(x,y)$ is flippable and by
$(x\,y)\circ\pi\circ(x\,y)$ otherwise.  $B=\{1\sim2\}$ (rows $1,2$ in a
common frame cycle), $A$ its complement; for $B$-instances $d_{12}$ and
$d_{21}$ are the numbers of $\pi$-steps from $1$ to $2$ and from $2$ to
$1$.

\begin{definition}[ladder]\label{def:ladder}
Put $K_0=\min(|C_1|,|C_2|)$ on $A$ and $K_0=\min(d_{12},d_{21})$ on
$B$, and for $1\le k<K_0$ let
\[
x_k=\pi^{-k}(1),\qquad y_k=\pi^{-k}(2),\qquad
T_k=\text{the $j'$-trade at }(x_k,y_k).
\]
The pairs $\{x_k,y_k\}$ are the \emph{ladder} of $(L,P)$, and
$\mathrm{Flip}(L,P)=\{k<K_0:(x_k,y_k)\text{ flippable}\}$.
\end{definition}

The rows $x_1,\dots,x_{K_0-1}$ are distinct rows of $C_1\setminus\{1\}$
(on $A$) or of the arc from $2$ to $1$ (on $B$), the $y_k$ likewise in
$C_2\setminus\{2\}$ or the arc from $1$ to $2$; so the $2(K_0-1)$
ladder rows are distinct and avoid rows $1,2$, and $K_0\ge2$ by the
admissibility of the mark ($\pi(1)\ne2$, $\pi(2)\ne1$).

\begin{theorem}[ladder identity]\label{thm:ladder}
For every $n$, $\lambda$ and split $\alpha+\beta=m$ with
$\alpha,\beta\ge2$,
\[
\Pr_X[B]-\Pr_X[A]\;=\;\Pr_X\bigl[B,\ \mathrm{Flip}=\emptyset\bigr]
-\Pr_X\bigl[A,\ \mathrm{Flip}=\emptyset\bigr],
\qquad\text{hence}\qquad
\bigl|\Pr_X[B]-\tfrac12\bigr|\;\le\;\tfrac12\Pr_X[\mathrm{Flip}=\emptyset].
\]
The same identity holds with $X$ replaced by $S(R)\times\{P\}$ for any
fixed admissible first-row pair $R$ and mark $P$.
\end{theorem}

\begin{proof}
Fix $(L,P)\in X$ and write $K=\{1,\dots,K_0-1\}$.  We show that the
$T_k$, $k\in K$, are commuting involutions of $X$ which preserve
$K_0$, the unordered ladder pairs and the set $\mathrm{Flip}$, and that
$T_k$ toggles $B$ exactly when $k\in\mathrm{Flip}$.  Granting this, the
group $G=\langle T_k:k\in K\rangle\cong(\mathbb Z/2)^{K}$ acts on $X$
with orbits $\{T_S(L,P):S\subseteq K\}$ of size $2^{|K|}$ (each $T_k$
changes rows $x_k,y_k$ and distinct $k$ change distinct rows), and on
the orbit of $(L,P)$ the bit of $T_S(L,P)$ is
$\beta(L,P)+|S\cap\mathrm{Flip}|\bmod2$.  If $\mathrm{Flip}\ne\emptyset$
exactly half of the subsets $S$ have $|S\cap\mathrm{Flip}|$ odd, so
$B$ and $A$ each contain half of the orbit; summing over orbits gives
the identity, and the bound follows from $\Pr[B]+\Pr[A]=1$.

\emph{$T_k$ is a well-defined involution of $X$.}  The $j'$-trade at
$(x_k,y_k)$ changes only rows $x_k,y_k\notin\{1,2\}$, so $\sigma$ and
the mark are untouched and $(L,P)$ stays in $X$; it is an involution
because the traded columns are again the $\rho'_{x_k,y_k}$-cycle
through $j'$ (\S\ref{sec:s13trapped}).

\emph{Effect on the frame.}  If $k\in\mathrm{Flip}$ the new frame is
$\pi'=(x_k\,y_k)\circ\pi$.  On $A$ this merges $C_1$ and $C_2$: the
new cycle runs $1\to\pi(1)\to\dots\to x_{k+1}\to y_k\to\pi(y_k)\to\dots
\to2\to\dots\to y_{k+1}\to x_k\to\dots\to1$, so $d'_{12}=|C_1|$ and
$d'_{21}=|C_2|$ (reading $x_{|C_1|}=1$, $y_{|C_2|}=2$).  On $B$,
$x_k$ lies on the arc from $2$ to $1$ and $y_k$ on the arc from $1$ to
$2$ (as $k<\min(d_{12},d_{21})$), and the split into
$(x_k,\pi(x_k),\dots,\pi^{-1}(y_k))$ and $(y_k,\dots,\pi^{-1}(x_k))$
puts $1$ ($k$ steps after $x_k$, fewer than the $d_{12}$ steps from
$x_k$ to $y_k$) in the first cycle and $2$ in the second: the new
state is in $A$ with $|C'_1|=d_{12}$, $|C'_2|=d_{21}$.  In both cases
$K_0$ is preserved and the bit toggles.  If $k\notin\mathrm{Flip}$ the
new frame is $(x_k\,y_k)\circ\pi\circ(x_k\,y_k)$: the cycle partition
is unchanged up to exchanging the labels $x_k,y_k$, so $K_0$ and the
bit are preserved.

\emph{The ladder is preserved.}  Let $\pi'=(x_k\,y_k)\circ\pi$ (the
conjugation case is similar and easier).  Since
$\pi'^{-1}(r)=\pi^{-1}((x_k\,y_k)(r))$, the preimage chains of $1$ and
$2$ under $\pi'$ agree with those under $\pi$ up to depth $k$ and are
exchanged beyond it: $\pi'^{-i}(1)=x_i$ and $\pi'^{-i}(2)=y_i$ for
$i\le k$, while $\pi'^{-(k+1)}(1)=\pi^{-1}(y_k)=y_{k+1}$ and then
$\pi'^{-i}(1)=y_i$, $\pi'^{-i}(2)=x_i$ for $k<i<K_0$.  So the unordered
pairs $\{x_i,y_i\}$, $i\in K$, are the ladder of $T_k(L,P)$.  In the
conjugation case the same computation with $(x_k\,y_k)\circ\pi\circ(x_k\,y_k)$
gives $\pi'^{-i}(1)=(x_k\,y_k)(x_i)$, which again permutes the pairs
$\{x_i,y_i\}$ among themselves (only $i=k$ is affected, by exchanging
$x_k,y_k$).

\emph{Flippability is preserved and the $T_k$ commute.}  $T_k$ changes
rows $x_k,y_k$ only, and the flippability of $(x_i,y_i)$, $i\ne k$,
depends only on rows $x_i,y_i$, which are different rows.  Likewise
$T_i$ and $T_k$ act on disjoint pairs of rows, each along a column set
determined by its own two rows, so $T_iT_k=T_kT_i$.  (After $T_k$ the
pair $\{x_i,y_i\}$ may have changed its index, but the trade at an
unordered pair does not depend on the index, so the group generated is
the same from every point of the orbit.)

The statement for fixed $R$ holds because every $T_k$ fixes rows $1$
and $2$.
\end{proof}

\begin{corollary}[the conjecture reduces to the ladder]\label{cor:ladder}
By Theorem~\ref{thm:marked}, $\mathrm{EQ}(\delta)$ holds with
$\delta=2\Pr_X[\mathrm{Flip}=\emptyset](1+o(1))$, and by
Theorem~\ref{thm:reduction}
\[
\dtv(\PP_n,\QQ_n)=O\Bigl(\log n\cdot\max_{\lambda,\alpha,\beta}
\Pr_{X(n,\lambda;\alpha,\beta)}[\mathrm{Flip}=\emptyset]+n^{-2/3}\Bigr).
\]
In particular the weak CGW conjecture follows from
\begin{equation}\label{eq:L}
\textup{(L)}\qquad \Pr_X\bigl[\text{no ladder pair is flippable}\bigr]=o(1/\log n)
\quad\text{uniformly in }\lambda,\alpha,\beta .
\end{equation}
\end{corollary}

\begin{remark}[what (L) asks, and what the data say]\label{rem:ladderdata}
(i) The ladder consists of $K_0-1$ pairs of \emph{disjoint} rows,
$K_0=\min(|C_1|,|C_2|)$ resp.\ $\min(d_{12},d_{21})$, and a pair is
flippable iff columns $j,j'$ lie in different cycles of its row
permutation --- the conjugate (rows$\leftrightarrow$columns) of the
bit itself.  (L) does not ask these coins to be fair, only that they
do not all fail.

(ii) Exact enumeration at $n=7$ (\texttt{s13\_comp7.c}, four
instances) confirms the identity to all digits
($\Pr[B]-\Pr[A]=-0.005848$ and $\Pr[B,\mathrm{Flip}=\emptyset]
-\Pr[A,\mathrm{Flip}=\emptyset]=-0.005848$ for $\lambda=(7)$,
$\alpha=2$; $+0.003509$ for $\alpha=3$; $-0.004651$ for $(5,2)$;
$0$ for a cross-cycle pair).  On uniform squares with true marks
(\texttt{s13\_ladder.py}, \texttt{runs/s13/f1/ladder\{30,50\}.log};
$9000$ and $7200$ instances):
\[
\Pr[\mathrm{Flip}=\emptyset\mid K_0]=2^{-(K_0-1)}
\]
to the resolution of the data for every $K_0$ ($n=50$: $0.480,0.265,
0.142,0.070,0.036,0.030,0.002$ for $K_0=2..8$ against $0.5,0.25,
0.125,0.0625,0.031,0.016,0.008$), the ladder coins are flippable with
frequency $0.4990$ ($n=30$) and $0.5002$ ($n=50$), and
$\Pr[\mathrm{Flip}=\emptyset]=0.125$ ($n=30$), $0.079$ ($n=50$)
$=\E[2^{-(K_0-1)}]$ to three digits.  Since $\Pr[K_0=k]\approx4/n$ for
small $k$, this is $\approx8/n$: if the coins are fair and
independent, (L) holds with $O(1/n)$ and the conjecture would follow
with $\dtv=O(\log n/n)$.

(iii) \emph{Cost of a short ladder.}  On $A$,
$\Pr[K_0\le\ell]\le72\ell/n$ by Proposition~\ref{prop:trapped}.  On
$B$ the ladder identity itself gives
$\Pr[B,\min(d_{12},d_{21})\le\ell,\mathrm{Flip}\ne\emptyset]
=\Pr[A,K_0\le\ell,\mathrm{Flip}\ne\emptyset]\le72\ell/n$, and the
switching of Proposition~\ref{prop:trapped} (splice a free cycle into
the short arc, toggling flippability by a turn of the free cycle)
gives $\Pr[B,d_{12}=\ell,|C_1|\le n/2]\le8/n$; what is missing is
$\Pr[B,\ d_{12}=\ell,\ |C_1|>n/2,\ \mathrm{Flip}=\emptyset]=O(1/n)$,
a short arc inside a cycle covering most rows with no flippable
ladder pair.  \emph{To do.}

(iv) \emph{Local witnesses.}  A ladder pair is flippable as soon as
the $\rho_{x_k,y_k}$-cycle through $j$ is short and avoids $j'$ (or
the one through $j'$ is short and avoids $j$).  On real squares the
length of the $j$-cycle of a ladder pair is spread like a uniform
variable on $[1,n]$ ($\Pr[\le\ell]=0.021,0.081,0.18,0.38$ for
$\ell=2,5,10,20$ at $n=50$), and among instances with $K_0\ge n/4$ a
flippable ladder pair with $j$-cycle of length $\le10$ exists in
$93\%$, $\le20$ in $99.6\%$.  So, granting (iii) with
$\ell_0=n/\log^3n$ (cost $O(1/\log^3n)$), (L) follows from
\begin{equation}\label{eq:Lloc}
\Pr\Bigl[K_0\ge\tfrac n{\log^3n}\ \text{and no ladder pair has a
$j$-cycle of length}\le\log^5n\text{ avoiding }j'\Bigr]=o(1/\log n),
\end{equation}
a statement about $\Theta(n/\log^3n)$ disjoint row pairs each of which
carries such a cycle with probability $\approx\log^5n/n$ (expected
number $\approx\log^2n$ witnesses).  Short cycles through a given
column in the row permutation of a given row pair are substructures
of polylogarithmic size, so \eqref{eq:Lloc} is the kind of statement
that the switching method of Kwan--Sudakov \cite{KS18} and its
successors is built for: a second-moment bound needs only
$\operatorname{Cov}(W_k,W_{k'})=o(\Pr[W_k]\Pr[W_{k'}]/\log^2n)$ for
disjoint pairs, in the space of completions of two rows and two
columns.  No fairness of any global statistic is needed anywhere.
\end{remark}

\begin{remark}[relation to the earlier programme]
Theorem~\ref{thm:ladder} is the orbit identity of the trade chain
(\S\ref{sec:s13tradechain}(b)) specialised to the abelian subgroup
generated by the ladder trades, for which the orbit structure is
explicit: on each orbit the bit is \emph{exactly} fair unless no ladder
pair is flippable.  Compared with the pool of intercalates
(Theorem~\ref{thm:master}) the ladder has deterministic bit-changing
pairs and an exact bias formula, at the price of needing the pairs
to be flippable; compared with the trade chain it needs no mixing.
The quenched version (fixed $R$) answers the fairness Question of
\texttt{paper/fairness\_question.tex}: $|2\Pr[\beta=1\mid R]-1|\le
\Pr[\mathrm{Flip}=\emptyset\mid R]$.  Items (I), (II), (III), (T),
(F1), (F2) and the component question are superseded by (L).
\end{remark}
